%% toc: P3 36 36.6
\chapter{Partially Ordered Quantifier Prefixes}
\label{sec:p03-w03-p03-044-po-prefixes}

In an ordinary first-order prefix each existentially quantified variable may
be chosen in the light of every universally quantified variable written to
its left.  Some statements need a weaker pattern of dependence: one witness
should be chosen from one universal value only, and a second witness from a
different universal value only.  No linear arrangement of \(\forall\) and
\(\exists\) expresses this.  A \emph{partially ordered} (or \emph{branching},
or \emph{Henkin}) prefix states the permitted dependences directly, and its
meaning is given by the existence of witness functions that consult only the
permitted arguments.

This chapter takes as input one-sorted structures with nonempty universes
and formulas of first-order logic or of first-order logic extended by the
Henkin quantifier.  Set-theoretic arguments are carried out in ZFC; the
Axiom of Choice is used exactly where it is named.  The conclusions are:
the general syntax of a dependency prefix and its function semantics; the
agreement of that semantics with ordinary satisfaction on linear prefixes;
the Henkin quantifier \(Q^{\mathrm H}\) and its two-row notation; a
sentence of \(\mathscr L_{\omega\omega}(Q^{\mathrm H})\) true in exactly the
infinite structures, and hence the failure of countable compactness; the
regularity of \(\mathscr L_{\omega\omega}(Q^{\mathrm H})\) on one-sorted
vocabularies, with the limit of its substitution property when infinitely
many constants are replaced; and the exact relation of \(Q^{\mathrm H}\) to the team
semantics of \cref{def:p03-team-semantics}.

Mundici's survey of further quantifiers in the same volume opens by turning
arbitrary partially ordered prefixes into quantifiers, which are
presented in \cref{sec:p03-w03-p03-042-mundici-po}; it continues with
quantifiers for comparing structures and quantifiers from Robinson
equivalence relations, and their treatment belongs with that section.

The Henkin quantifier and the infinity sentence follow
Ebbinghaus's presentation in \emph{Model-Theoretic Logics}
\cite[Ch.~II, Section~2.5, pp.~38--39]{Ebbinghaus1985ExtendedLogics}.
The general dependency-triple syntax follows Zdanowski
\cite[Section~2, pp.~2--3]{Zdanowski2016HenkinArxiv}; that presentation is
the arXiv version~1 of 21 December 2016.  The two sources use different
letters, and each keeps its own: Zdanowski writes a structure \(M\) and a
prefix \((A,E,D)\), while Ebbinghaus writes a structure \(\mathfrak A\) whose
universe is \(A\).  In
\cref{def:p03-w03-p03-044-henkin-prefix,def:p03-w03-p03-044-linear-triple,thm:p03-w03-p03-044-linear-skolem}
the letter \(A\) names a set of universal variables; from
\cref{def:p03-w03-p03-044-qh} on it names the universe of \(\mathfrak A\).

\section{Dependency prefixes and their function semantics}
\label{sec:p03-w03-p03-044-prefix-syntax}

\begin{hbdefinition}{def:p03-w03-p03-044-henkin-prefix}{Henkin prefix as a dependency triple}
A \emph{Henkin prefix} is an ordered triple \(Q=(A,E,D)\) in which \(A\)
and \(E\) are disjoint finite sets of variables, the \emph{universal} and
the \emph{existential} variables of \(Q\), and \(D\subseteq A\times E\) is
the \emph{dependency relation}.  A variable \(y\in E\) \emph{depends on}
\(x\in A\) when \((x,y)\in D\)
\cite[p.~2]{Zdanowski2016HenkinArxiv}.  Zdanowski does not distinguish a
quantifier from its prefix, and neither does this chapter.

Write \(A=\{x_1,\ldots,x_n\}\) and \(E=\{y_1,\ldots,y_k\}\) with the
displayed indexing fixed.  For each \(i\le k\), let \(\overline{x}_i\) be
the list, in increasing order of index, of the variables \(x_j\) with
\((x_j,y_i)\in D\), and let \(m_i\) be its length.  Let \(\varphi\) be a
formula, \(M\) a structure with nonempty universe and \(s\) an assignment
defined on the free variables of \(\varphi\) outside \(A\cup E\).  Then
\(M\entails Q\varphi(x_1,\ldots,x_n,y_1,\ldots,y_k)\,[s]\) holds iff there
are functions \(f_i:M^{m_i}\to M\) \((i\le k)\) such that for all
\(a_1,\ldots,a_n\in M\),
\[
 M\entails\varphi\bigl[a_1,\ldots,a_n,
   f_1(\overline a_1),\ldots,f_k(\overline a_k)\bigr]\,[s],
\]
where \(\overline a_i\) is the subtuple of \((a_1,\ldots,a_n)\) at the
positions of \(\overline x_i\).  When \(m_i=0\), a function
\(f_i:M^{0}\to M\) is identified with the element it returns.  Zdanowski
states the same clause as truth of
\(\forall\overline{\mathbf x}\,\varphi(\overline{\mathbf x},f_1(\overline
x_1),\ldots,f_k(\overline x_k))\) in the expansion
\((M,\{f_i\}_{i\le k})\), where \(\overline{\mathbf x}\) lists all universal
variables of \(Q\) \cite[pp.~2--3]{Zdanowski2016HenkinArxiv}.  The variables
of \(A\cup E\) are bound by \(Q\); the free variables of \(Q\varphi\) are
those of \(\varphi\) outside \(A\cup E\).

For a family \(\mathcal Q\) of Henkin prefixes, \(L(\mathcal Q)\) is
first-order logic extended by the quantifiers in \(\mathcal Q\), and
\(L(Q)\) abbreviates \(L(\{Q\})\); \(\mathcal H\) denotes the family of all
Henkin prefixes \cite[p.~3]{Zdanowski2016HenkinArxiv}.

The name \emph{partially ordered} is exact.  Because \(A\cap E=\varnothing\),
no pair \((v,v)\) lies in \(D\), and no two pairs of \(D\) can be chained,
since the second entry of a pair lies in \(E\) and the first entry of a pair
lies in \(A\).  Hence \(D\) is irreflexive and (vacuously) transitive, so it
is a strict partial order on \(A\cup E\).

The witness functions are the intended strategies.  Read \(Q\varphi\) as a
game in which a falsifier supplies all universal values and a verifier
supplies each existential value \(y_i\) after seeing only the universal
values \(\overline a_i\) on which \(y_i\) depends.  A verifier strategy is
then exactly a tuple \((f_1,\ldots,f_k)\) as above, and it wins when the
matrix holds for every choice of universal values.  The clause above says
that such a winning strategy exists.
\index{Henkin quantifier}\index{partially ordered quantifier prefix}
\index{branching quantifier}
\end{hbdefinition}

\begin{hbdefinition}{def:p03-w03-p03-044-linear-triple}{The dependency triple of a linear prefix}
Let \(\pi=\mathsf Q_1z_1\cdots\mathsf Q_nz_n\) be an ordinary prefix of
distinct variables with each \(\mathsf Q_i\in\{\forall,\exists\}\).  Its
\emph{dependency triple} is
\(Q_\pi=(A_\pi,E_\pi,D_\pi)\), where
\(A_\pi=\{z_i:\mathsf Q_i=\forall\}\), \(E_\pi=\{z_j:\mathsf Q_j=\exists\}\),
and
\[
 D_\pi=\{(z_i,z_j): i<j,\ \mathsf Q_i=\forall,\ \mathsf Q_j=\exists\}.
\]
Variables in \(A_\pi\) and \(E_\pi\) are indexed in order of occurrence in
\(\pi\).  Thus every existential variable depends on every universal variable
written before it.  Ebbinghaus calls this the \emph{linear character} of a
usual first-order prefix and notes that Skolem functions make it visible
\cite[p.~38]{Ebbinghaus1985ExtendedLogics}.
\end{hbdefinition}

\begin{hbtheorem}{thm:p03-w03-p03-044-linear-skolem}{proved=annex:p03-w03-p03-044-linear-skolem}{Linear prefixes and witness functions}
Let \(\pi\) be a linear prefix as in
\cref{def:p03-w03-p03-044-linear-triple}, let \(\varphi\) be a formula of
\(\mathscr L_{\omega\omega}\) or of \(\mathscr L_{\omega\omega}(Q^{\mathrm
H})\), let \(M\) be a structure with nonempty universe and let \(s\) be an
assignment defined on the free variables of \(\pi\varphi\).
\begin{enumerate}
 \item If \(M\entails Q_\pi\varphi\,[s]\) in the sense of
       \cref{def:p03-w03-p03-044-henkin-prefix}, then \(M\entails\pi\varphi\,[s]\)
       in the ordinary sense.
 \item Assume the Axiom of Choice.  If \(M\entails\pi\varphi\,[s]\), then
       \(M\entails Q_\pi\varphi\,[s]\).
\end{enumerate}
\end{hbtheorem}

The proof uses only the clauses for \(\forall\) and \(\exists\); the matrix
\(\varphi\) is treated through its satisfaction relation alone.  Choice
enters in the universal step of part~(2), where one witness family must be
selected for every value of the universal variable at once.

\begin{hbexample}{def:p03-w03-p03-044-linear-example}{Skolem functions for a linear prefix}
Ebbinghaus's illustration is the prefix
\(\forall u\,\exists v\,\forall wx\,\exists y\)
\cite[p.~38]{Ebbinghaus1985ExtendedLogics}.  Its dependency triple is
\[
 \bigl(\{u,w,x\},\ \{v,y\},\ \{(u,v),(u,y),(w,y),(x,y)\}\bigr).
\]
So \(v\) receives a unary witness function and \(y\) a ternary one, and by
\cref{thm:p03-w03-p03-044-linear-skolem}
\(\forall u\,\exists v\,\forall wx\,\exists y\,\varphi(u,v,w,x,y)\) is
equivalent, under the Axiom of Choice, to the existence of \(f\) and \(g\)
with
\[
 \forall uwx\,\varphi\bigl(u,f(u),w,x,g(u,w,x)\bigr),
\]
which Ebbinghaus writes with function quantifiers as
\[
 \exists fg\,\forall uwx\,\varphi\bigl(u,f(u),w,x,g(u,w,x)\bigr)
\]
\cite[p.~38]{Ebbinghaus1985ExtendedLogics}.
\end{hbexample}

\section{The Henkin quantifier}
\label{sec:p03-w03-p03-044-henkin-quantifier}

\begin{hbdefinition}{def:p03-w03-p03-044-qh}{The Henkin quantifier and the Henkin logic}
The formulas of \(\mathscr L_{\omega\omega}(Q^{\mathrm H})\) are generated
by the first-order formation rules together with one more rule: if
\(\varphi\) is a formula and \(x_0,y_0,x_1,y_1\) are distinct variables, then
\(Q^{\mathrm H}x_0y_0x_1y_1\,\varphi\) is a formula, in which the four
displayed variables are bound.  Let \(\mathfrak A\) be a structure with
universe \(A\).  Ebbinghaus's clause is
\[
 \mathfrak A\entails Q^{\mathrm H}x_0y_0x_1y_1\,
   \varphi(x_0,y_0,x_1,y_1)
\]
iff there are functions \(f_0,f_1:A\to A\) such that for all
\(a_0,a_1\in A\),
\(\mathfrak A\entails\varphi[a_0,f_0(a_0),a_1,f_1(a_1)]\)
\cite[p.~38]{Ebbinghaus1985ExtendedLogics}.  Further free variables of
\(\varphi\) are held fixed by an assignment, as in
\cref{def:p03-w03-p03-044-henkin-prefix}.  The quantifier binds four
variables and is the simplest prefix that is not of linear character
\cite[p.~38]{Ebbinghaus1985ExtendedLogics}.

To display the dependences, the same formula is written as
\[
 \begin{matrix}\forall x_0\,\exists y_0\\ \forall x_1\,\exists y_1\end{matrix}
 \ \varphi(x_0,y_0,x_1,y_1)
\]
\cite[p.~39]{Ebbinghaus1985ExtendedLogics}.  Each row is read separately:
\(y_0\) is chosen from \(x_0\) alone, and \(y_1\) from \(x_1\) alone.  In
the notation of \cref{def:p03-w03-p03-044-henkin-prefix}, \(Q^{\mathrm H}\)
is the triple
\[
 \bigl(\{x_0,x_1\},\ \{y_0,y_1\},\ \{(x_0,y_0),(x_1,y_1)\}\bigr),
\]
and the clause above is the general clause for this triple: the witness
for \(y_0\) has argument list \((x_0)\), and the witness for \(y_1\) has
argument list \((x_1)\).  Zdanowski writes the same prefix as
\[
 \bigl(\{x,z\},\ \{y,w\},\ \{(x,y),(z,w)\}\bigr)
\]
and calls it \(\mathsf H\) \cite[p.~2]{Zdanowski2016HenkinArxiv}.  Ebbinghaus calls
\(\mathscr L_{\omega\omega}(Q^{\mathrm H})\) the \emph{Henkin logic} and
attributes the quantifier to Henkin \cite[p.~38]{Ebbinghaus1985ExtendedLogics}.
\end{hbdefinition}

\begin{hbtheorem}{thm:p03-w03-p03-044-qh-vs-linear}{proved=annex:p03-w03-p03-044-qh-vs-linear}{Removing a dependence strengthens the prefix}
Let \(\varphi\) be a formula of \(\mathscr L_{\omega\omega}\) or of
\(\mathscr L_{\omega\omega}(Q^{\mathrm H})\), \(\mathfrak A\) a structure and
\(s\) an assignment for the remaining free variables of \(\varphi\).
\begin{enumerate}
 \item If \(\mathfrak A\entails Q^{\mathrm H}x_0y_0x_1y_1\,\varphi\,[s]\),
       then \(\mathfrak A\entails\forall x_0\,\exists y_0\,\forall x_1\,
       \exists y_1\,\varphi\,[s]\).
 \item The converse fails: in the empty vocabulary, with
       \(A=\{0,1\}\) and \(\varphi\) the formula \(y_1=x_0\), the linear
       sentence holds and the Henkin sentence fails.
\end{enumerate}
\end{hbtheorem}

The linear prefix \(\forall x_0\,\exists y_0\,\forall x_1\,\exists y_1\)
has the dependency relation
\[
 \{(x_0,y_0),\ (x_0,y_1),\ (x_1,y_1)\},
\]
and the Henkin prefix deletes the pair \((x_0,y_1)\).  A verifier who may consult fewer values has
fewer strategies, so the Henkin sentence is the stronger one.  Part~(2) shows
that the loss can be decisive.

\begin{hbtheorem}{thm:p03-w03-p03-044-qh-infinity}{proved=annex:p03-w03-p03-044-qh-infinity}{The Henkin quantifier defines infinity}
Let \(\varphi_\infty\) be the sentence of the empty vocabulary
\[
 \exists z\
 \begin{matrix}\forall x_0\,\exists y_0\\ \forall x_1\,\exists y_1\end{matrix}
 \bigl(z\neq y_0\land(y_0=x_1\to y_1=x_0)\bigr)
\]
\cite[p.~39]{Ebbinghaus1985ExtendedLogics}.  Let \(\mathfrak A\) be a
structure of any vocabulary with universe \(A\neq\varnothing\).
\begin{enumerate}
 \item \(\mathfrak A\entails\varphi_\infty\) iff there are \(a\in A\) and
       \(f_0,f_1:A\to A\) with \(a\notin\operatorname{rg}(f_0)\) and
       \(f_1(f_0(b))=b\) for all \(b\in A\).
 \item If \(\mathfrak A\entails\varphi_\infty\), then \(A\) is infinite.
 \item Assume the Axiom of Choice.  If \(A\) is infinite, then
       \(\mathfrak A\entails\varphi_\infty\).
\end{enumerate}
\end{hbtheorem}

Ebbinghaus summarizes the three parts by saying that the condition in~(1)
means that \(A\) is infinite \cite[p.~39]{Ebbinghaus1985ExtendedLogics}; the
direction from infinity to the condition is where choice is used.  The
condition says that \(A\) admits a one-to-one map into itself that misses a
point, the witness \(f_1\) undoing \(f_0\) on its range.

\begin{hbexample}{def:p03-w03-p03-044-qh-example}{Evaluating the infinity sentence}
Let \(\mathfrak A=(\N)\) in the empty vocabulary.  Take \(z=0\),
\(f_0(n)=n+1\), and \(f_1(0)=0\), \(f_1(n)=n-1\) for \(n\ge1\).  For any
\(a_0,a_1\), the value \(y_0=a_0+1\) differs from \(0\), and if
\(a_0+1=a_1\) then \(y_1=f_1(a_1)=a_0\).  So the verifier wins and
\(\mathfrak A\entails\varphi_\infty\).

Let \(A=\{0,1,2\}\).  Any \(f_0\) with a left inverse \(f_1\) is one-to-one
on a three-element set and therefore onto, so no \(z\) lies outside its
range.  Hence \(A\) does not satisfy \(\varphi_\infty\): whatever pair of
row strategies the verifier fixes, some choice of \(x_0,x_1\) defeats it.
The first conjunct \(z\neq y_0\) forces \(f_0\) to miss \(z\), and the
second forces \(f_1\) to undo \(f_0\).
\end{hbexample}

\begin{hbcorollary}{thm:p03-w03-p03-044-qh-not-compact}{proved=annex:p03-w03-p03-044-qh-not-compact}{\(\mathscr L_{\omega\omega}(Q^{\mathrm H})\) is not \(\aleph_0\)-compact}
Assume the Axiom of Choice.  There is a countable set of
\(\mathscr L_{\omega\omega}(Q^{\mathrm H})\)-sentences of the empty
vocabulary each finite subset of which has a model and which has no model.
Thus \(\mathscr L_{\omega\omega}(Q^{\mathrm H})\) is not
\(\aleph_0\)-compact \cite[p.~39]{Ebbinghaus1985ExtendedLogics}.
\end{hbcorollary}

Here a logic \(\mathscr L\) is \(\kappa\)-compact, for an infinite cardinal
\(\kappa\), when every set \(\Phi\) of \(\mathscr L\)-sentences of one
vocabulary, of power at most \(\kappa\), has a model as soon as each finite
subset of \(\Phi\) has one
\cite[Ch.~II, Definition~1.2.4(i), p.~31]{Ebbinghaus1985ExtendedLogics}.
\(\aleph_0\)-compactness is the countable case.

\begin{hbcorollary}{thm:p03-w03-p03-044-qh-not-fo}{proved=annex:p03-w03-p03-044-qh-not-fo}{No first-order equivalent of \(\varphi_\infty\)}
Assume the Axiom of Choice.  No first-order sentence \(\theta\) of the empty
vocabulary has the same nonempty models as \(\varphi_\infty\).
\end{hbcorollary}

The proof applies first-order compactness directly to the negation of
such a \(\theta\) together with the sentences requiring arbitrarily
large finite cardinalities. Thus
\(\mathscr L_{\omega\omega}(Q^{\mathrm H})\) is strictly more expressive than
\(\mathscr L_{\omega\omega}\), and the extra expressive power costs
countable compactness, which first-order logic has by
\cref{thm:fol-compactness}.

\section{The Henkin logic is regular}
\label{sec:p03-w03-p03-044-regular}

Ebbinghaus states, directly after the two-row notation, that the Henkin logic
\(\mathscr L_{\omega\omega}(Q^{\mathrm H})\) is regular
\cite[p.~39]{Ebbinghaus1985ExtendedLogics}, in the sense of the closure
properties of \cref{def:p03-w03-p03-047-closure,def:p03-w03-p03-047-relativization}
(his Definitions~1.2.1--1.2.3) for a logic in the sense of
\cref{def:p03-w03-p03-047-ebbinghaus-logic}.  This section makes the logic
precise as an assignment of sentence sets to vocabularies and proves the
claim.  The clause for \(Q^{\mathrm H}\) is written for structures with a
single universe \(A\), with witness functions \(f_0,f_1:A\to A\)
\cite[p.~38]{Ebbinghaus1985ExtendedLogics}, and the relativization and
substitution properties are formulated only for one-sorted vocabularies
\cite[Ch.~II, p.~30]{Ebbinghaus1985ExtendedLogics}.  Accordingly the
vocabularies below are one-sorted.

\begin{hbconvention}{def:p03-w03-p03-044-henkin-syntax}{Formulas, assignments and sentence sets of the Henkin logic}
Let \(\tau\) be a one-sorted vocabulary of relation symbols, function
symbols and constants.  The \(\tau\)-terms are built from variables and
constants by the function symbols.  The formulas of
\(\mathscr L_{\omega\omega}(Q^{\mathrm H})\) of vocabulary \(\tau\) are the
atomic formulas \(t_1=t_2\) and \(Rt_1\cdots t_n\), and \(\neg\theta\),
\((\theta_0\land\theta_1)\), \(\exists x\,\theta\) and
\(Q^{\mathrm H}x_0y_0x_1y_1\,\theta\) for distinct \(x_0,y_0,x_1,y_1\);
\(\forall\), \(\lor\), \(\to\) and \(\leftrightarrow\) abbreviate their
usual definitions.  In this section an assignment in \(\mathfrak A\) is a
total map \(s\) from the variables into \(A\), and \(s[a/x]\) changes the
value at \(x\) to \(a\).  The clauses are the first-order ones together
with the following: \(\mathfrak A\entails Q^{\mathrm
H}x_0y_0x_1y_1\,\theta\,[s]\) iff there are \(f_0,f_1:A\to A\) such that
\(\mathfrak A\entails\theta\,[s^{f_0f_1}_{a_0a_1}]\) for all
\(a_0,a_1\in A\), where \(s^{f_0f_1}_{a_0a_1}=s[a_0/x_0][f_0(a_0)/y_0][a_1/x_1][f_1(a_1)/y_1]\);
this is the clause of \cref{def:p03-w03-p03-044-qh} with the remaining free
variables read off \(s\).  A \emph{sentence} is a formula without free
variables.  \(\mathscr L_{\omega\omega}(Q^{\mathrm H})[\tau]\) is the set of
sentences all of whose symbols lie in \(\tau\), and for
\(\varphi\in\mathscr L_{\omega\omega}(Q^{\mathrm H})[\tau_{\mathfrak A}]\),
\(\mathfrak A\entails\varphi\) means \(\mathfrak A\entails\varphi\,[s]\) for
some assignment \(s\), equivalently, by
\cref{thm:p03-w03-p03-044-basic}(1), for every \(s\).

A formula is \emph{flat} if each of its atomic subformulas has one of the
forms \(v=w\), \(Rv_1\cdots v_n\), \(fv_1\cdots v_n=w\) or \(c=w\), with
\(v,w,v_1,\ldots,v_n\) variables.  For an injective map \(\iota\) from
variables to variables, \(\iota\theta\) replaces every occurrence of every
variable \(v\) in \(\theta\), bound ones included, by \(\iota(v)\).  For
distinct constants \(c_1,\ldots,c_m\) and variables \(v_1,\ldots,v_m\), not
necessarily distinct, \(\theta(\bar v/\bar c)\) replaces each occurrence of
\(c_i\) by \(v_i\).  For a renaming \(\rho:\tau\to\sigma\) in the sense of
\cref{def:p03-w03-p03-047-vocabularies}, \(\theta^\rho\) replaces each symbol
\(S\) of \(\tau\) by \(\rho(S)\).
\end{hbconvention}

\begin{hblemma}{thm:p03-w03-p03-044-basic}{proved=annex:p03-w03-p03-044-basic}{Basic properties of satisfaction}
Let \(\theta\) be a formula of \(\mathscr L_{\omega\omega}(Q^{\mathrm H})\)
of the one-sorted vocabulary \(\tau\).
\begin{enumerate}
 \item \emph{Coincidence.}  If \(\mathfrak A\) and \(\mathfrak B\) have the
       same universe and interpret every symbol occurring in \(\theta\)
       alike, and \(s,s'\) agree on the free variables of \(\theta\), then
       \(\mathfrak A\entails\theta\,[s]\) iff
       \(\mathfrak B\entails\theta\,[s']\).
 \item \emph{Isomorphism.}  If \(\pi:\mathfrak A\cong\mathfrak B\) for
       \(\tau\)-structures, then \(\mathfrak A\entails\theta\,[s]\) iff
       \(\mathfrak B\entails\theta\,[\pi\circ s]\).
 \item \emph{Symbol renaming.}  If \(\rho:\tau\to\sigma\) is a renaming,
       then \(\mathfrak A\entails\theta\,[s]\) iff
       \(\mathfrak A^\rho\entails\theta^\rho\,[s]\).
 \item \emph{Variable renaming.}  If \(\iota\) is an injective map from
       variables to variables, then \(\mathfrak A\entails\iota\theta\,[s]\)
       iff \(\mathfrak A\entails\theta\,[s\circ\iota]\).
 \item \emph{Constants to variables.}  Let \(c_1,\ldots,c_m\in\tau\) be
       distinct constants, \(\tau_0=\tau\setminus\{c_1,\ldots,c_m\}\), and
       \(v_1,\ldots,v_m\) variables none of which occurs in \(\theta\).  If
       \(s(v_i)=c_i^{\mathfrak A}\) for every \(i\), then
       \(\mathfrak A\entails\theta\,[s]\) iff
       \(\mathfrak A\restriction\tau_0\entails\theta(\bar v/\bar c)\,[s]\).
\end{enumerate}
\end{hblemma}

\begin{hblemma}{thm:p03-w03-p03-044-unnesting}{proved=annex:p03-w03-p03-044-unnesting}{Every formula has a flat equivalent}
For every formula \(\theta\) of \(\mathscr L_{\omega\omega}(Q^{\mathrm H})\)
of vocabulary \(\tau\) and every finite set \(W\) of variables there is a
flat formula \(\theta^{\mathrm u}\) of vocabulary \(\tau\) with the same
free variables, every variable of which either occurs in \(\theta\) or lies
outside \(W\), such that \(\mathfrak A\entails\theta\,[s]\) iff
\(\mathfrak A\entails\theta^{\mathrm u}\,[s]\) for every \(\tau\)-structure
\(\mathfrak A\) and every assignment \(s\).
\end{hblemma}

\begin{hbtheorem}{thm:p03-w03-p03-044-regular}{proved=annex:p03-w03-p03-044-regular}{Regularity of the Henkin logic}
On one-sorted vocabularies, \(\mathscr L_{\omega\omega}(Q^{\mathrm H})\)
with the sentence sets and satisfaction of
\cref{def:p03-w03-p03-044-henkin-syntax} has the following properties.
\begin{enumerate}
 \item It satisfies (i)--(v) of \cref{def:p03-w03-p03-047-ebbinghaus-logic},
       with \(\varphi^\rho\) as the renamed sentence.
 \item It has the atom, negation, conjunction and particularization
       properties of \cref{def:p03-w03-p03-047-closure}; for a constant
       \(c\in\tau\) and \(\varphi\in\mathscr L_{\omega\omega}(Q^{\mathrm
       H})[\tau]\) one may take \(\exists x\,\varphi(x/c)\) with \(x\) a
       variable not occurring in \(\varphi\).
 \item It has the relativization property of
       \cref{def:p03-w03-p03-047-relativization}.
 \item Substitution: let \(\tau\subseteq\tau'\), let
       \(\varphi\in\mathscr L_{\omega\omega}(Q^{\mathrm H})[\tau']\), and for
       each symbol \(S\in\tau'\setminus\tau\) let
       \(\psi_S(\mathbf c_S)\in\mathscr L_{\omega\omega}(Q^{\mathrm
       H})[\tau\cup\{\mathbf c_S\}]\) be given as in
       \cref{def:p03-w03-p03-047-relativization}.  Let \(S_1,\ldots,S_k\) be
       the symbols of \(\tau'\setminus\tau\) that occur in \(\varphi\).  There
       is \(\varphi^\ast\in\mathscr L_{\omega\omega}(Q^{\mathrm H})[\tau]\)
       such that a \(\tau\)-structure \(\mathfrak A\) satisfies
       \(\varphi^\ast\) iff \(\mathfrak A\) has an expansion \(\mathfrak A'\)
       to \(\tau'\) with \(\mathfrak A'\entails\varphi\) in which, for
       \(j\le k\), the relation \(S_j^{\mathfrak A'}\) is the set defined by
       \(\psi_{S_j}\) and the function or constant \(S_j^{\mathfrak A'}\) has
       as graph the set defined by \(\psi_{S_j}\).  If only
       finitely many symbols of \(\tau'\setminus\tau\) are function symbols
       or constants (in particular, if \(\tau'\setminus\tau\) is finite), the
       same holds with the condition imposed on every symbol of
       \(\tau'\setminus\tau\).
\end{enumerate}
Hence \(\mathscr L_{\omega\omega}(Q^{\mathrm H})\), on one-sorted
vocabularies, is a regular logic.  No use of the Axiom of Choice is made.
\end{hbtheorem}

The proof treats \(Q^{\mathrm H}\) exactly like a first-order quantifier in
each induction: its clause quantifies over two functions on the universe,
and each construction transports, extends or restricts those two functions.
In the relativization the matrix becomes
\((\chi_{x_0}\land\chi_{x_1})\to(\chi_{y_0}\land\chi_{y_1}\land\theta^U)\),
where \(\chi_v\) expresses membership of \(v\) in the relativizing set; a
witness pair on the large structure restricts to the relativized universe
because that universe is nonempty, and a witness pair on the relativized
universe extends arbitrarily outside it.  For the substitution property the
sentence is first made flat (\cref{thm:p03-w03-p03-044-unnesting}), so that
each replaced symbol occurs only in an atom with variable arguments, and each
such atom is then replaced by the corresponding \(\psi_S\) with its constants
turned into those variables.

The restriction in part~(4) concerns symbols that do not occur in
\(\varphi\).  Ebbinghaus displays the meaning of a single replacement of a
constant \(c\) by \(\lambda d\,\psi(d)\) as
\(\exists c(\forall d(c=d\leftrightarrow\psi(d))\land\varphi)\)
\cite[Ch.~II, p.~30]{Ebbinghaus1985ExtendedLogics}, and this sentence is
false when \(\psi\) does not define a single element, whether or not \(c\)
occurs in \(\varphi\).  A relation symbol imposes no such condition, since
every definable set is a relation.  With infinitely many function symbols or
constants the conditions cannot in general be collected into one sentence,
even in first-order logic:

\begin{hbtheorem}{thm:p03-w03-p03-044-substitution-limit}{proved=annex:p03-w03-p03-044-substitution-limit}{Substitution of infinitely many constants}
Assume that the set of variables is countable.  Let \(\tau=\{e\}\) with
\(e\) a constant, and for \(n\ge2\) let \(\varepsilon_n\) be the
first-order sentence
\(\exists v_1\cdots\exists v_n(\bigwedge_{1\le i<j\le n}v_i\neq v_j\land
\forall v\,\bigvee_{i=1}^n v=v_i)\), which holds in a structure iff its
universe has exactly \(n\) elements.  There is a set
\(X\subseteq\{n\in\mathbb N:n\ge2\}\) with the following property.  Let
\(\tau'=\tau\cup\{c_n:n\in X\}\) with distinct new constants \(c_n\), let
\(\varphi\) be the sentence \(e=e\), and for \(n\in X\) let
\(\psi_{c_n}(d)\) be the first-order sentence \(d=e\land\neg\varepsilon_n\)
of vocabulary \(\tau\cup\{d\}\).  Then no sentence \(\varphi^\ast\) of
\(\mathscr L_{\omega\omega}(Q^{\mathrm H})[\tau]\) satisfies, for every
\(\tau\)-structure \(\mathfrak A\):
\(\mathfrak A\entails\varphi^\ast\) iff \(\mathfrak A\) has an expansion
\(\mathfrak A'\) to \(\tau'\) with \(\mathfrak A'\entails\varphi\) in which,
for every \(n\in X\), \(\{c_n^{\mathfrak A'}\}\) is the set defined by
\(\psi_{c_n}\).  In particular no first-order sentence of vocabulary
\(\tau\) has this property.
\end{hbtheorem}

The proof shows that the required class consists of the
\(\tau\)-structures whose size is not in \(X\), and that there are more such
classes than sentences.  The data of the theorem are first-order, so the
same failure occurs for \(\mathscr L_{\omega\omega}\) itself.  Part~(4) of
\cref{thm:p03-w03-p03-044-regular} therefore gives the substitution property
in the two forms that are available for first-order logic as well: for the
symbols that occur in \(\varphi\), and for every symbol when only finitely
many function symbols and constants are replaced.

The theorem settles the first claim on p.~39.  The text there continues that
adjoining to \(\mathscr L_{\omega\omega}\) quantifiers like \(Q^{\mathrm H}\)
that arise from partially ordered prefixes yields the full expressive power
of second-order logic, and refers to Section~VI.1 of the same volume for the
details \cite[p.~39]{Ebbinghaus1985ExtendedLogics}.  This chapter neither
states that claim as a theorem nor proves it: the passage on p.~39 does not
fix which family of prefixes is adjoined or in what sense the expressive
power is meant.  A precise result of that section is Mundici's theorem that
every class projective in first-order logic is elementary in the logic of
all partially ordered prefixes (\cref{thm:p03-w03-p03-042-mundici-pc-ec}).

\section{Henkin quantification and team semantics}
\label{sec:p03-w03-p03-044-teams}

The dependence logic of \cref{sec:p03-dependence-logics} expresses
restricted dependence by an atom \(=(\bar x,y)\) evaluated on a team, not by
a new quantifier.  The two devices meet exactly when the matrix is
first-order, through the following flatness property.

\begin{hblemma}{thm:p03-w03-p03-044-flatness}{proved=annex:p03-w03-p03-044-flatness}{First-order formulas are flat on teams}
Let \(\psi\) be a formula of the dependence-logic grammar of
\cref{def:p03-team-semantics} containing no dependence atom.  For every
structure \(M\) with nonempty universe and every team \(X\) on \(M\) whose
variable domain contains the free variables of \(\psi\),
\[
 M,X\entails\psi\iff(\forall s\in X)\ M,s\entails\psi,
\]
where the right side is ordinary Tarskian satisfaction.
\end{hblemma}

\begin{hbtheorem}{thm:p03-w03-p03-044-qh-team}{proved=annex:p03-w03-p03-044-qh-team}{The Henkin quantifier as a dependence-logic sentence}
Assume the Axiom of Choice.  Let \(\psi\) be a formula of the
dependence-logic grammar without dependence atoms, let \(x_0,y_0,x_1,y_1\)
be distinct variables, let \(\mathfrak A\) be a structure with nonempty
universe \(A\), and let \(s\) be an assignment defined on the free variables
of \(\psi\) outside \(\{x_0,y_0,x_1,y_1\}\).  Then
\[
 \mathfrak A\entails Q^{\mathrm H}x_0y_0x_1y_1\,\psi\,[s]
 \iff
 \mathfrak A,\{s\}\entails
   \forall x_0\,\exists y_0\,\forall x_1\,\exists y_1
   \bigl(=(x_1,y_1)\land\psi\bigr).
\]
\end{hbtheorem}

The atom \(=(x_1,y_1)\) removes the dependence of \(y_1\) on \(x_0\) and
\(y_0\) that the linear prefix would otherwise allow.  No atom is needed for
\(y_0\), since only \(x_0\) precedes it.  For the matrix of
\(\varphi_\infty\), the implication is first rewritten with literals as
\(z\neq y_0\land(y_0\neq x_1\lor y_1=x_0)\), so that it lies in the
dependence grammar.

\begin{hbwarning}{def:p03-w03-p03-044-if-distinct}{Slashed quantifiers are a different convention}
The independence-friendly quantifier \((\exists x/W)\) of
\cref{def:p03-if-logic} records hidden variables on one existential move of
a game.  The Henkin quantifier is instead one four-place quantifier symbol
added to first-order logic, with the function clause of
\cref{def:p03-w03-p03-044-qh}.  Both express that a choice ignores certain
earlier values, but this chapter proves a translation only for dependence
atoms with a first-order matrix
(\cref{thm:p03-w03-p03-044-qh-team}).  It does not identify
\(\mathscr L_{\omega\omega}(Q^{\mathrm H})\) with independence-friendly logic
or with dependence logic.  In \(\mathscr L_{\omega\omega}(Q^{\mathrm H})\),
\(Q^{\mathrm H}\) may occur under negation and inside its own matrix, while
dependence logic has no negated dependence atoms.
\end{hbwarning}

\begin{hbwarning}{def:p03-w03-p03-044-limits}{Choice, empty universes and the cost in compactness}
The equivalence of a linear prefix with its witness functions, the
definition of infinity by \(\varphi_\infty\), and the team translation each
use the Axiom of Choice in one direction.  Without it, part~(2) of
\cref{thm:p03-w03-p03-044-linear-skolem} and part~(3) of
\cref{thm:p03-w03-p03-044-qh-infinity} are not available from the proofs
given here.  All clauses assume a nonempty universe, because a witness
function of arity zero must return an element.  Finally, adding even the
single quantifier \(Q^{\mathrm H}\) to first-order logic already destroys
countable compactness (\cref{thm:p03-w03-p03-044-qh-not-compact}).
Compactness arguments for first-order logic therefore do not carry over to
\(\mathscr L_{\omega\omega}(Q^{\mathrm H})\), although the closure
properties of a regular logic do
(\cref{thm:p03-w03-p03-044-regular}).  That theorem concerns one-sorted
vocabularies only, since the clause for \(Q^{\mathrm H}\) used here is
written for a single universe.
\end{hbwarning}

\paragraph{Sources.}
Henkin quantifier, two-row notation, linear-prefix illustration and
infinity sentence: Ebbinghaus
\cite[Ch.~II, Section~2.5, pp.~38--39]{Ebbinghaus1985ExtendedLogics}.
Compactness notion:
\cite[Ch.~II, Definition~1.2.4, p.~31]{Ebbinghaus1985ExtendedLogics}.
Regularity claim: \cite[p.~39]{Ebbinghaus1985ExtendedLogics}; structures,
\(\tau\)-closed sets, renamings, the notion of logic and the regularity
properties:
\cite[Ch.~II, \S1.1, pp.~27--28, and Definitions~1.2.1--1.2.3, pp.~29--31]{Ebbinghaus1985ExtendedLogics},
as set out in \cref{def:p03-w03-p03-047-vocabularies,def:p03-w03-p03-047-ebbinghaus-logic,def:p03-w03-p03-047-closure,def:p03-w03-p03-047-relativization}.
General dependency-triple syntax and function semantics: Zdanowski
\cite[Section~2, pp.~2--3]{Zdanowski2016HenkinArxiv}.  Team semantics and
the independence-friendly quantifier: \cref{sec:p03-dependence-logics}.
